% P10, 15.09.2026: Gate-Reichweite und Host-Zustimmung; StewardGateTools/RunTools/ChatRunner.
% M25, 13.09.2026: bestehende Dokumentationsbelege in den numerischen Zitierstil überführt.
% M10, 11.09.2026: Leserführung, Kapiteltrennung und belegte Reichweite.
% StewardChatRunner.cs (Agent, Session, Reducer); StewardGateTools.cs:368--397;
% IngestGateDecisions.cs:24--45 (gleicher Loader, zwei schema-gleiche Dateien);
% StewardRunTools.cs (Approval, RunCoreAnalysis, Resume); HandoffSpikeTests.cs.
% B-32: gemeinsame Einreichung ist hier ein Vertrag, keine identische Schreibfunktion.

\section{Steward-Realisierung}
\label{sec:realisierung-steward}

Der Steward aus
Abschnitt~\ref{sec:durable-ausfuehrung-agentische-bedienung} ist ein
einzelner MAF-Agent mit Systemauftrag, Werkzeugkatalog und
Dialogschleife. Er nutzt damit die Agentabstraktion für die
Interpretation von Anliegen und die Werkzeugwahl. Die aufgerufenen
Workflows behalten die Kontrolle über ihre Verarbeitungsschritte.

MAF ermöglicht die Serialisierung einer Sitzung, sodass ihr
Gesprächszusammenhang gespeichert und später wieder geladen werden
kann \cite{msmaf2026session}.
Das Artefakt legt diese Sitzungen dateibasiert ab. Für den an das
Modell übergebenen Verlauf sind drei Strategien konfigurierbar:
vollständige Nachrichtenfolge, begrenztes Nachrichtenfenster oder
zusammenfassende Verdichtung. Diese Verlaufsreduktion wird über die
entsprechende Frameworkschnittstelle in die Dialogführung eingebunden
\cite{msmaf2026storage}.
Die Sitzung dient der Fortführung des Dialogs; maßgebliche
Projektinformationen werden weiterhin aus dem Core bezogen.

Der Werkzeugkatalog erschließt vorhandene Systemfähigkeiten:
Projekt- und Laufstatus lesen, Verarbeitung starten oder wiederaufnehmen
und Gate-Entscheidungen einreichen. Ein Teil der Lesezugriffe nutzt die
MCP-Anbindung aus Abschnitt~\ref{sec:realisierung-tools-mcp}.
Auch die Core-Analyse wird über denselben Aufruf wie im Direktzugang
gestartet. Für entsprechend ausgewiesene kosten- oder
außenwirkungsträchtige Werkzeuge greift die dort beschriebene
Zustimmungsprüfung vor der Ausführung.

Gate-Entscheidungen lassen sich für ausgewählte Gate-Arten direkt im
Dialog bearbeiten, etwa am Ingest-Gate. Für andere Entscheidungspunkte
kann der Steward eine unterstützte Review-Oberfläche öffnen oder auf
den gesonderten Review-Aufruf verweisen. Vor der Ausführung eines
Einreichungswerkzeugs zeigt die Dialogsteuerung den konkreten Aufruf
an und fragt die Zustimmung des Menschen ab.

Der Übergang vom Dialog zur fachlichen Wirkung lässt sich am
Diktateingang erklären: Der Mensch äußert einen Präzisierungswunsch;
der Steward fasst ihn als Eingabe und startet nach Zustimmung den
zuständigen Workflow. Dessen Änderungsvorschlag wird am Ingest-Gate
vorgelegt. Der Steward fragt die Entscheidung ab und überträgt sie mit
dem zustimmungspflichtigen Einreichungswerkzeug in eine
Entscheidungsdatei. Geprüft werden dabei die betroffenen Kennungen,
zulässige Optionen, vollständige Einträge und bei Ablehnung eine
Begründung. Für diesen Eingang verwenden Steward und Review-Oberfläche
Dateien mit demselben Schema, die über dieselbe Ladefunktion in die
Gate-Antwort eingehen. Die Einreichung kann unmittelbar die
Wiederaufnahme anstoßen. Ein bloß besprochener Vorschlag ändert den
Core noch nicht; ohne eingereichte Antwort bleibt der Lauf pausiert.
Fall F2 in Abschnitt~\ref{sec:eval-resume} dokumentiert eine solche
Aufgabe, ohne damit jede Dialogvariante zu prüfen.

Die Kontrolle verteilt sich auf drei Ebenen: Der Systemauftrag
beschreibt das erwartete Dialogverhalten; die Werkzeugdefinitionen
begrenzen die angebotenen Fähigkeiten; Werkzeug-, Workflow- und
Speicherlogik prüfen die jeweiligen Aufrufe und Zustandsänderungen.
Die Aufforderung an den Agenten, nachzufragen, ist daher von einer
technischen Prüfung zu unterscheiden. Der Werkzeugkatalog bietet
keinen unmittelbaren Core-Schreibzugriff. Eine umfassende Untersuchung
möglicher Umgehungswege oder der Bedienqualität erfolgte nicht.

Die Handoff-Orchestrierung wurde in einem isolierten Test mit
skriptgesteuerten Modellantworten erprobt, aber nicht für den Steward
übernommen. Sie übergibt die Aufgabenbearbeitung an einen anderen
Agenten \cite{msmaf2026handoff}; beim gewählten Werkzeugmodell bleibt dagegen
der Steward der dialogführende Agent. Der Test belegt die technische
Übergabe, keine Überlegenheit einer der beiden Steuerungsformen.
Zusätzliche Dienste für dauerhaftes Agentengedächtnis wurden nicht
eingebunden; die Sitzungskontinuität wird durch die beschriebene
Speicherung gedeckt. Wie die fachlichen Wirkungen der aufgerufenen
Abläufe gespeichert und nach außen übertragen werden, behandelt der
folgende Abschnitt.
